\section{Motivación}

% Introduccion al problema
La navegación autónoma es un área de investigación activa en robótica, con todo tipo de aplicaciones, desde vehículos autónomos hasta robots de servicio y drones. La capacidad de un robot para navegar de forma segura y eficiente en entornos desconocidos es fundamental para su utilidad y adopción en entornos reales. Esto se logra mediante la construcción de mapas del entorno y la localización del robot dentro de esos mapas, lo que se conoce como \textit{Simultaneous Localization and Mapping} (SLAM). 

A lo largo de la literatura, se han propuesto numerosos algoritmos de SLAM, cada uno con sus propias ventajas y desventajas. Dentro de esos algoritmos, hay un especial interés en aquellos basados en visión, ya que las cámaras son sensores económicos y proporcionan una gran cantidad de información sobre el entorno. Sin embargo, los algoritmos que las utilizan son computacionalmente costosos, lo que limita su aplicación en tareas como la navegacion autónoma, que requieren respuestas a tiempo real.

Una forma de clasificar estos algoritmos es según dos factores: El primero, con qué sensores trabajan, y el segundo, qué tipo de representación del entorno utilizan. Los sensores con los que cuenta el robot determinan la información con la que trabaja el algoritmo, el cual deberá fusionarla entre sí para que brinden resultados más robustos, en un marco de coherencia con la representación del entorno.

% Sensores
En este trabajo se consideran robots equipados con sensores visuales e inerciales, una configuración conocida como VINS (\textit{Visual-Inertial SLAM}). En particular, se empleará una cámara a color junto con un mapa de profundidad por píxel, y una IMU (unidad de medición inercial) que aporta datos de aceleraciones y velocidades angulares en los tres ejes. La información visual permite observar la estructura y la apariencia del entorno, mientras que la IMU entrega medidas de la dinámica del movimiento entre las imágenes, información esencial para mantener estimaciones de pose robustas ante cambios rápidos y en áreas poco texturizadas. La elección de estos sensores se debe a su bajo costo y alta disponibilidad, al igual que a su factibilidad, pues estos sensores son sencillos de integrar y calibrar. Además, su reducido tamaño, requerimiento energético y peso los hacen ideales para plataformas móviles con restricciones de carga útil, como drones o robots de servicio. Sin embargo, la información que aportan es ruidosa y a menudo insuficiente, brindando resultados poco robustos en comparación a la integración de sensores como LiDARs, los cuales son más precisos pero traen consigo un costo mayor en energía, peso y valor monetario.

% Modelo del mundo
Como modelo de mapa, este trabajo adopta la técnica de Gaussian Splatting, surgida recientemente \cite{kerbl2023gaussian}, que ha despertado gran interés por su combinación de calidad y eficiencia de renderizado. En este modelo, el entorno se representa como una colección de gaussianas tridimensionales, cada una con su propia posición, forma y color. Estas gaussianas pueden ser renderizadas de manera eficiente utilizando técnicas de rasterización, lo que permite obtener imágenes de alta calidad a partir de la representación del entorno. Además, el modelo de Gaussian Splatting es capaz de representar detalles finos y texturas complejas, por lo que resulta adecuado para aplicaciones de navegación autónoma en entornos reales. En este modelo, existen distintas técnicas para inicializar y optimizar las gaussianas, surgiendo mucha literatura al respecto.

% VIGS-Fusion / Loop Closure
Dentro del marco de SLAM en el que se ubica este trabajo, artículos recientes exploran distintas formas de integrar los sensores con el modelo de mundo. Tal es el caso de VIGS-Fusion \cite{ndoye2025vigs}, que presenta un algoritmo de SLAM eficiente en tiempo real. Sin embargo, una propiedad que no posee es la que se conoce en la literatura como la detección y cierre de ciclos, que es la capacidad del algoritmo de reconocer lugares previamente visitados y corregir la trayectoria en consecuencia. Del mismo modo, al encontrarse la técnica de Gaussian Splatting en etapas muy tempranas de desarrollo, existe muy poca exploración al respecto, consistiendo mayormente en artículos dispersos orientados a combinaciones específicas de sensores, sin suficiente trabajo en la integración de los resultados entre ellos.

